% =============================================================================
% Utility Commands
% =============================================================================



% =============================================================================
% =============================================================================
% Convex Hull Path
% =============================================================================

\newcommand{\convexpath}[2]{
[   
    create hullnodes/.code={
        \global\edef\namelist{#1}
        \foreach [count=\counter] \nodename in \namelist {
            \global\edef\numberofnodes{\counter}
            \node at (\nodename) [draw=none,name=hullnode\counter] {};
        }
        \node at (hullnode\numberofnodes) [name=hullnode0,draw=none] {};
        \pgfmathtruncatemacro\lastnumber{\numberofnodes+1}
        \node at (hullnode1) [name=hullnode\lastnumber,draw=none] {};
    },
    create hullnodes
]
($(hullnode1)!#2!-90:(hullnode0)$)
\foreach [
    evaluate=\currentnode as \previousnode using \currentnode-1,
    evaluate=\currentnode as \nextnode using \currentnode+1
    ] \currentnode in {1,...,\numberofnodes} {
  let
    \p1 = ($(hullnode\currentnode)!#2!-90:(hullnode\previousnode)$),
    \p2 = ($(hullnode\currentnode)!#2!90:(hullnode\nextnode)$),
    \p3 = ($(\p1) - (hullnode\currentnode)$),
    \n1 = {atan2(\y3,\x3)},
    \p4 = ($(\p2) - (hullnode\currentnode)$),
    \n2 = {atan2(\y4,\x4)},
    \n{delta} = {-Mod(\n1-\n2,360)}
  in 
    {-- (\p1) arc[start angle=\n1, delta angle=\n{delta}, radius=#2] -- (\p2)}
}
-- cycle
}


% =============================================================================
% Reflect Node Utility
% =============================================================================

\newcommand{\reflectnode}[1]{%
\draw[red, thick] (#1.center) circle (0.11);
\fill[red, opacity=0.3] (#1.center) circle (0.11);
\node[red, line width=8, ultra thick, scale=1.5] at (#1.center)
  {\Huge $\times$};
}

% =============================================================================
% KB Circle (from tikzmacros.tex)
% =============================================================================

\def\KBcircle{\tikz[baseline=.1ex]{%
    \draw[color=gray,thick] (-0.3,0.3) -- (0.3,-0.3);
    \draw[color=gray,thick] (-0.3,-0.3) -- (-0.05,-0.05);
    \draw[color=gray,thick] (0.05,0.05) -- (0.3,0.3);
}}

% =============================================================================
% Quiver Curve Style
% =============================================================================

\tikzset{curve/.style={settings={#1},to path={(\tikztostart)
    .. controls ($(\tikztostart)!\pv{pos}!(\tikztotarget)!\pv{height}!270:(\tikztotarget)$)
    and ($(\tikztostart)!1-\pv{pos}!(\tikztotarget)!\pv{height}!270:(\tikztotarget)$)
    .. (\tikztotarget)\tikztonodes}},
    settings/.code={\tikzset{quiver/.cd,#1}
        \def\pv##1{\pgfkeysvalueof{/tikz/quiver/##1}}},
    quiver/.cd,pos/.initial=0.35,height/.initial=0}

% =============================================================================
% FreeTikZ Thickness (shared default line width)
% =============================================================================

\def\thickness{0.7pt}

% =============================================================================
% On-the-fly TikZ SVG compiler and includer
% =============================================================================
\RequirePackage{svg}

\newcommand{\ctikzsvg}[2][width=\linewidth]{%
  \begingroup
  \CatchFileEdef\abspath{|kpsewhich "#2"}{\endlinechar=-1 }%
  \immediate\write18{python3 /home/dzack/.pandoc/bin/render_figures.py \abspath}%
  \filename@parse{#2}%
  \begin{center}%
    \includesvg[#1]{/home/dzack/.pandoc/figures/rendered/\filename@base}%
  \end{center}%
  \endgroup
}

